% Discussion --- defensible claims and honest limitations. This section
% is the guard against over-claiming; it states explicitly what we do NOT
% claim.
\section{Discussion}
\label{sec:discussion}

\subsection{Defensible claims}
Our results support four claims, stated conservatively.
\begin{enumerate}[leftmargin=*]
    \item Temporal transformers with a phenological branch are the
    strongest individual perceivers on PASTIS-R (TSViT-pheno, mIoU
    0.6253), and heterogeneous stacking over decorrelated members
    lifts the held-out parcel-level macro F1 to 0.7470
    (accuracy 0.8490), $+12.3$ points over the best single model. The
    same champion, used as the agent's perceiver, raises parcel accuracy
    from 0.8308 to 0.9413 ($+11$ points) on 13{,}481 real parcels,
    confirming that the Be My Eyes separation pays off precisely because
    the frozen reasoner inherits the champion's accuracy.
    \item A contrastive vision-language branch~\cite{li2025farslip}
    contributes a useful but modest complementary signal (the
    five-member stacking gain over three members is $+0.0118$ macro
    F1). It is a complement, not a standalone classifier.
    \item The \emph{Be My Eyes} pattern~\cite{huang2025bemyeyes} yields
    a grounded conversational layer in which every cited figure is
    traceable to an auditable tool result, with a frozen reasoner that
    is interchangeable between cloud (Gemini~2.5~Pro) and on-premises
    (Qwen3.5-35B-A3B) for data sovereignty.
    \item Spatial transfer from France to Catalonia is not zero-shot;
    few-shot fine-tuning recovers mIoU 0.2468, in line with reported
    transferability limits~\cite{harvesting2026alphaearth}.
\end{enumerate}

\subsection{What we explicitly do not claim}
To avoid over-claiming, we delimit the boundaries of our evidence.
\begin{itemize}[leftmargin=*]
    \item We do \emph{not} claim that a fine-tuned VLM outperforms the
    frozen Gemini reasoner. The language layer is a communication,
    explanation, and tool-orchestration component, not a pixel
    classifier; its value is grounding, auditability, and the
    on-premises sovereignty option.
    \item We do \emph{not} claim a classification threshold
    (e.g.\ F1~$\geq$~0.80) for Mexico. The avocado/guava demonstration
    over Mexican parcels is a qualitative, zero-shot methodological
    demo without ground truth; it illustrates the pipeline, not a
    measured accuracy.
    \item We do \emph{not} claim that zero-shot transfer works outside
    France. The France$\rightarrow$Catalonia experiment shows zero-shot
    failure and few-shot recovery; a $k$-shot curve on
    EuroCropsML~\cite{reuss2025eurocropsml} characterises the data
    budget required.
    \item We do \emph{not} claim TSViT mIoU above 0.75. PASTIS-R
    saturates near 0.70 mIoU for this class set, and the phenological
    branch contributes on the order of $0.3$ points, not several.
\end{itemize}

\subsection{Limitations}
Three structural limitations remain. First, six minority classes with
ambiguous phenology drag the global macro F1 down; addressing them
requires additional data or targeted augmentation, not more models of
the same type. Second, cloud cover can obscure a large fraction of
pixels in rainy seasons, motivating fusion with Sentinel-1 radar.
Third, the conversational-layer benchmark on
AgroMind~\cite{ruan2025agromind} is pending the H100 evaluation
(US-069); we report the protocol and column structure but no
synthetic numbers.

\subsection{Cost efficiency (FinOps)}
The system is engineered to run at low operational cost. The four
serving components (API, frontend, tiling, inference worker) run on
Cloud Run with scale-to-zero, and the on-premises H100 reasoner is an
Azure spot instance activated only on demand (roughly three hours per
day). The steady-state operational cost is approximately
US\$115 per month, while the one-time training budget --- the H100 and
L4 spot windows used to fit the segmenters and the FarSLIP distillation
--- totals US\$262--602. The frozen-reasoner design is itself a cost
lever: because the LLM is never fine-tuned as a classifier, no
expensive GPU fine-tuning of a frontier model is on the critical path,
and the heavy compute is confined to the one-time perception-stack
training.

\subsection{Reproducibility}
Every figure in Section~\ref{sec:results} is anchored to a versioned
artefact via \texttt{\% src:} comments in the LaTeX source; datasets and
weights are tracked with DVC and experiment runs with MLflow under
\texttt{data\_version} and \texttt{code\_version} tags. The repository
is released under Apache~2.0.
